% ThermalQualityRevision20261023-v1
\hypersetup{bookmarksdepth=3}
\begingroup
\renewcommand{\Needspace}[1]{\par\begingroup
\dimen0=#1\relax\vskip0pt plus\dimen0\penalty-100
\vskip0pt plus-\dimen0\vskip\dimen0\penalty9999\vskip-\dimen0
\vskip0pt\endgroup}
\tcbset{breakable/.default=false}
\chapter{有心力、万有引力与天体运动}
% Editable TikZ: use the force centre, not the centre of an ellipse, for distances.
\providecommand{\ThermalFigure}[3]{\begin{qualityfigure}
\begin{tikzpicture}[>=Latex,font=\small,x=.8cm,y=.8cm]#1\end{tikzpicture}
\figcap{#2}\label{#3}\end{qualityfigure}}
\newcommand{\ThRadialFallBase}[1]{
 \draw[->](-.5,0)--(7,0)node[right]{$r$};\fill(0,0)circle(4pt)node[below]{$M$};
 \fill[DeepBlue](2.5,0)circle(2pt)node[above]{$r_1$};\fill[DeepBlue](6,0)circle(2pt)node[above]{$r_0$};
 \draw[->,thick](5.5,.7)--(3.9,.7)node[midway,above]{下落};#1}
\newcommand{\ThQRadialFall}{\ThermalFigure{\ThRadialFallBase{}}{从外侧静止释放，径向坐标随时间减小，始末距离都为正。}{ch11:quality:radial-q}}
\newcommand{\ThARadialFall}{\ThermalFigure{\ThRadialFallBase{
 \draw[<->](3,-.8)--node[below]{$|\dd r|$}(3.7,-.8);\node[below]at(4.5,-1.6){$\dot r<0,\quad\dd t=-\dd r/|\dot r|$};
}}{径向速度取负分支，时间元为正；积分上下限须与此相配。}{ch11:quality:radial-a}
图~\ref{ch11:quality:radial-a}中位移箭头朝向力心，故能量方程求出的速率须配 $\dot r=-\sqrt{2GM(1/r-1/r_0)}$。时间积分若交换为从 $r_1$ 到 $r_0$，负号才随之消失；这一符号核对避免得到负的下落时间。
}
\newcommand{\ThOrbitBase}[1]{
 \fill[DeepBlue!12](0,0)circle(1.4);\draw[thick](0,0)circle(1.4);\node at(0,0){地球};
 \draw[dashed](0,0)circle(1.65);\fill[DeepBlue](1.65,0)circle(2pt)node[right]{$P$};
 \draw[->](0,0)--(140:1.4)node[midway,above]{$R$};#1}
\newcommand{\ThQOrbit}{\ThermalFigure{\ThOrbitBase{
 \node[right]at(3,1){近地轨道：高度相对 $R$ 可忽略};
}}{圆轨道速度和逃逸速度对应不同的运动要求；图中间隙为示意。}{ch11:quality:orbit-q}}
\newcommand{\ThAOrbit}{\ThermalFigure{\ThOrbitBase{
 \draw[->,thick](1.65,0)--(1.65,1.5)node[right]{$v_c$};
 \draw[->,DeepBlue,thick](1.65,0)--(.5,0)node[above]{$F_g$};
 \node[right,align=left]at(3,0){圆轨道：引力提供向心力\\逃逸临界：无穷远速度趋零};
}}{用 $GM=gR^2$ 把引力参数换为题给地表量；两速度不能混用。}{ch11:quality:orbit-a}
图~\ref{ch11:quality:orbit-a}中的圆周条件是 $mv_c^2/R=GMm/R^2$。逃逸临界态则以无穷远为势能零点并令总机械能为零；二者分别由力与能量决定，故逃逸速度比圆轨道速度大 $\sqrt2$ 倍。
}
\newcommand{\ThApsidesBase}[1]{
 \draw[thick](0,0)ellipse(3 and 1.8);\fill(2.4,0)circle(3pt)node[above left]{$S$};
 \fill(3,0)circle(2pt)node[right]{$P$};\fill(-3,0)circle(2pt)node[left]{$A$};
 \draw[<->](2.4,-2.2)--node[below]{$r_p$}(3,-2.2);\draw[<->](-3,-3)--node[below]{$r_a$}(2.4,-3);#1}
\newcommand{\ThQApsides}{\ThermalFigure{\ThApsidesBase{}}{近点 $P$、远点 $A$ 的距离从焦点力心 $S$ 量起；椭圆中心不作为距离原点。}{ch11:quality:apsides-q}}
\newcommand{\ThTransferBase}[1]{
 \draw[TextGray](0,0)circle(1.2);\draw[TextGray](0,0)circle(3.2);
 \draw[DeepBlue,thick](-1,0)ellipse(2.2 and 1.96);\fill(0,0)circle(3pt)node[below right]{$S$};
 \fill(1.2,0)circle(2pt)node[below right]{$P$};\fill(-3.2,0)circle(2pt)node[left]{$A$};
 \draw[->](0,0)--(50:1.2)node[right]{$r_1$};\draw[->](0,0)--(120:3.2)node[above]{$r_2$};#1}
\newcommand{\ThQTransfer}{\ThermalFigure{\ThTransferBase{}}{转移椭圆与内、外圆轨道分别相切于 $P,A$；力心为它们共同的焦点 $S$。}{ch11:quality:transfer-q}}
\newcommand{\ThATransfer}{\ThermalFigure{\ThTransferBase{
 \draw[->,thick](1.2,0)--(1.2,1.3)node[right]{$\Delta v_1$};
 \draw[->,thick](-3.2,0)--(-3.2,-1)node[left]{$\Delta v_2$};
 \draw[->,DeepBlue,thick](.35,1.55)--(-.55,1.9);
}}{两次脉冲均沿当地顺行切向；无推力段从近点飞到远点，历时半个椭圆周期。}{ch11:quality:transfer-a}
图~\ref{ch11:quality:transfer-a}中的第一脉冲把内圆速度提高到转移椭圆近点速度；第二脉冲把较低的远点速度提高到外圆速度。两处箭头在纸面方向相反，却都沿顺行方向，因此两次速度增量的大小均为正。
}
\newcommand{\ThMassLossBase}[1]{
 \draw[dashed](-1.2,0)circle(1.2);\draw[dashed](-1.2,0)circle(3);
 \fill[DeepBlue](-2.4,0)circle(5pt)node[below]{$M$};\fill(1.8,0)circle(3pt)node[below]{$m$};
 \fill(-1.2,0)circle(1.5pt)node[below]{$C$};\draw[<->](-2.4,-1)--node[below]{$d$}(1.8,-1);
 \draw[->,thick](-2.4,0)--(-2.4,-1.5);\draw[->,thick](1.8,0)--(1.8,1.5);#1}
\newcommand{\ThQMassLoss}{\ThermalFigure{\ThMassLossBase{
 \foreach\a in {30,90,150,210,270,330}{\draw[->,TextGray](-2.4,0)--++(\a:.8);}
 \node[right]at(3,0){$M$ 瞬时各向同性抛出 $\Delta M$};
}}{爆发前绕共同质心 $C$ 转动；无反冲使残余两星瞬时速度连续。}{ch11:quality:mass-loss-q}}
\newcommand{\ThAMassLoss}{\ThermalFigure{
 \fill[DeepBlue](0,0)circle(4pt)node[below]{$M-\Delta M$};\fill(5,0)circle(3pt)node[below]{$m$};
 \draw[<->](0,-.8)--node[below]{$d$（瞬时不变）}(5,-.8);
 \draw[->,thick](5,0)--(5,1.8)node[right]{$v_{\rm rel}$（瞬时不变）};
 \node[above]at(2.5,2.4){相对运动的引力参数改变为 $G(M+m-\Delta M)$};
}{爆发后转入残余双星的相对坐标，束缚条件由相对比能量为负判断。}{ch11:quality:mass-loss-a}
图~\ref{ch11:quality:mass-loss-a}固定的是爆发瞬间的相对距离与相对速度；爆发前的质心位置和分量动能不应继续当作残余系统的质心量。用 $v_{\rm rel}^2=G(M+m)/d$ 与新的引力参数比较，才能直接判断新轨道是否束缚。
}
\newcommand{\ThLRLBase}[1]{
 \fill(0,0)circle(3pt)node[below]{$O$};\draw[->](0,0)--(5,0)node[right]{极轴};
 \fill[DeepBlue](2.4,1.7)circle(2pt)node[right]{$P$};\draw[->](0,0)--(2.4,1.7)node[midway,above left]{$\vect r$};
 \draw(.8,0)arc(0:35.3:.8)node[right]{$\theta$};#1}
\newcommand{\ThQLRL}{\ThermalFigure{\ThLRLBase{
 \draw[->,thick](2.4,1.7)--(1.3,2.5)node[above]{$\vect p$};
}}{平面有心力运动中的位置、动量与极角；待证明守恒矢量可用于选取极轴。}{ch11:quality:lrl-q}}
\newcommand{\ThALRL}{\ThermalFigure{\ThLRLBase{
 \draw[->,DeepBlue,thick](0,0)--(2,0)node[below]{$\vect e$};
 \draw[->,thick](0,0)--(4.8,3.4)node[right]{$\uvec r$ 的方向};
}}{将固定的偏心率矢量作为极轴，点乘径向单位矢量引入 $e\cos\theta$。}{ch11:quality:lrl-a}
图~\ref{ch11:quality:lrl-a}中的 $\theta$ 从常矢量 $\vect e$ 量起，故 $\vect e\cdot\uvec r=e\cos\theta$。另一方面，$\vect L$ 垂直轨道平面；利用三重积把 $(\vect p\times\vect L)\cdot\uvec r$ 化为 $L^2/r$，便把矢量守恒式转成轨道方程。圆轨道 $e=0$ 时极轴可任取。
}
\newcommand{\ThQDust}{\ThermalFigure{
 \draw[thick](0,0)circle(.7);\node at(0,0){$M_s$};\draw[->](0,0)--(-45:.7)node[below]{$r_s$};
 \draw[->,DeepBlue,thick](-4,1.8)--(-1.8,1.8)node[midway,above]{$v$};\draw[dashed](-4,0)--(0,0);
 \draw[<->](-3.7,0)--node[left]{$b$}(-3.7,1.8);\node[right]at(2,1.6){远处平行来流，密度 $\rho$};
}{飞船静止系中的入射尘埃；冲量参数 $b$ 从远处渐近线量起，不等于最近距离。}{ch11:quality:dust-q}}
\newcommand{\ThQPrecession}{\ThermalFigure{
 \draw[thick](0,0)circle(2);\fill(0,0)circle(3pt)node[below]{$M$};\fill[DeepBlue](2,0)circle(2pt)node[above]{$m$};
 \draw[->](2,0)--(.8,0)node[above]{$F_g$};\draw[->,DeepBlue](2,0)--(3,0);
 \node[DeepBlue]at(3.6,.7){$F_1=kmR$};\draw[thin,TextGray](3,0)--(3.4,.45);
 \draw[<->](0,-.7)--node[below]{$R$}(2,-.7);\node[align=left,anchor=west]at(5,1.3){近圆轨道上\\存在微小径向扰动};
}{径向附加力朝外；圆周角运动与径向小振动须分开求频率。}{ch11:quality:precession-q}}
\newcommand{\ThAPrecession}{\ThermalFigure{
 \draw[TextGray,domain=0:360,samples=100]plot({2*cos(\x)/(1+.18*cos(\x))},{2*sin(\x)/(1+.18*cos(\x))});
 \draw[DeepBlue,thick,domain=0:360,samples=100]plot({2*cos(\x+25)/(1+.18*cos(\x))},{2*sin(\x+25)/(1+.18*cos(\x))});
 \fill(0,0)circle(2pt)node[below]{$O$};\draw[dashed](0,0)--(1.695,0);\draw[dashed](0,0)--(25:1.695);
 \draw[->](1,0)arc(0:25:1);\node[right]at(2.8,.5){$\Delta\varpi=2\pi(\Omega/\kappa_r-1)$};
}{两次近日点方向的示意叠图，不把缓慢进动轨道误认为一条闭合椭圆。}{ch11:quality:precession-a}
图~\ref{ch11:quality:precession-a}中一次近日点到下一次近日点用的是径向周期 $2\pi/\kappa_r$，这期间方位角增加 $2\pi\Omega/\kappa_r$。超出 $2\pi$ 的部分就是顺行进动角；因此进动角速度取两频率之差，而不是把附加力简单代入开普勒周期。
}
\newcommand{\ThQOrbitDrag}{\ThermalFigure{
 \draw[thick](0,0)ellipse(3 and 1.8);\fill(2.4,0)circle(3pt)node[below]{$M$};\fill[DeepBlue](0,1.8)circle(2pt)node[above]{$m$};
 \draw[->,thick](0,1.8)--(-1.4,1.8)node[above]{$\vect v$};\draw[->,DeepBlue,thick](0,1.8)--(1.4,1.8)node[above]{$-\gamma\vect v$};
}{弱阻力始终与瞬时速度反向；研究对象是很多周期后的缓慢平均演化。}{ch11:quality:drag-q}}
\newcommand{\ThAOrbitDrag}{\ThermalFigure{
 \draw[thick](-1.8,0)ellipse(3 and 2.4);\draw[DeepBlue,thick](-1.08,0)ellipse(1.8 and 1.44);
 \fill(0,0)circle(3pt)node[below]{$S$};\draw[->](-4,1.6)--(-2.7,1.05);
 \node[right,align=left]at(2,1){同一焦点 $S$\\周期平均下 $a$ 减小、$e$ 保持不变};
}{以力心为中心按比例缩放的相继密切椭圆；实际轨道不在一个周期内突然跳变。}{ch11:quality:drag-a}
图~\ref{ch11:quality:drag-a}中的两条椭圆共用焦点 $S$，两个半轴按同一因子缩小。证明这一点需要同时追踪 $E$ 与 $L$：由它们的平均变化检验 $e^2=1+2EL^2/(G^2M^2m^3)$ 不变，不能只由能量降低就断言偏心率不变。
}



第9章已建立角动量守恒、有效势与有心力轨道的基本方法。本章先作必要回顾，再重点讨论轨道初始条件、二体运动及天体运动中的具体应用。

\section{有心力运动的守恒量}
\label{ch11:sec:central}

把力心固定在惯性系原点 $O$。若质点所受合力始终沿 $O$ 与质点的连线，就称它为有心力。
先研究大小只依赖距离、且不显含时间的情形：
\begin{equation}\label{ch11:eq:central-force}
 \vect F=F(r)\uvec r,\qquad r=|\vect r|>0.
\end{equation}
这里 $F(r)$ 是有符号的径向分量：$F>0$ 向外，$F<0$ 向内。
“沿径向”保证力矩为零；“只依赖 $r$、不显含时间”才使这里的机械能也守恒。
例如随时间改变的径向力仍可使角动量守恒，却未必有守恒的机械能。
真实的两颗星都会运动，暂时固定力心是质量悬殊时的近似；后面将用相对坐标解除这个限制。

对原点的角动量为 $\vect L=m\vect r\times\vect v$。由于叉乘两个平行矢量为零，
\begin{equation}
 \frac{\dd\vect L}{\dd t}=\vect r\times\vect F=0.
\end{equation}
若 $\vect L\ne0$，任意时刻都有 $\vect r\cdot\vect L=0$，所以轨道在一个通过力心、垂直于
$\vect L$ 的固定平面内。若 $L=0$，速度与位置矢量平行，运动退化为沿直线的径向运动；
这一情形不能套用后面含 $1/L$ 或 $1/h$ 的轨道公式。

在非退化轨道平面内取极坐标 $(r,\theta)$，选择 $\theta$ 的正向使角动量大小 $L>0$，记
\begin{equation}\label{ch11:eq:h}
 h=\frac Lm=r^2\dot\theta>0.
\end{equation}
$h$ 是单位质量角动量，量纲为 $\mathrm{m^2/s}$。极坐标的两个单位矢量也随时间转动，故
\[
 \dot{\uvec r}=\dot\theta\uvec\theta,\qquad
 \dot{\uvec\theta}=-\dot\theta\uvec r,
\]
\[
 \vect v=\dot r\uvec r+r\dot\theta\uvec\theta,\qquad
 \vect a=(\ddot r-r\dot\theta^2)\uvec r
 +(r\ddot\theta+2\dot r\dot\theta)\uvec\theta.
\]
牛顿第二定律的径向、切向分量分别是
\begin{equation}\label{ch11:eq:polar-newton}
 m(\ddot r-r\dot\theta^2)=F(r),\qquad
 r\ddot\theta+2\dot r\dot\theta=0.
\end{equation}
将后一式乘以 $r$，左端正是 $\dd(r^2\dot\theta)/\dd t$，再次得到式\eqref{ch11:eq:h}。
在短时间 $\dd t$ 内位置矢量扫过面积 $\dd A=\tfrac12r^2\dd\theta$，因此
\begin{equation}\label{ch11:eq:area-law}
 \frac{\dd A}{\dd t}=\frac12r^2\dot\theta=\frac h2.
\end{equation}
这就是开普勒（Kepler）第二定律的动力学来源；等面积定律适用于有心力，并非平方反比引力独有。

再取势能 $U(r)$ 满足 $F(r)=-\dd U/\dd r$。两个速度分量正交，机械能为
\begin{equation}\label{ch11:eq:energy-original}
 E=\frac12m(\dot r^2+r^2\dot\theta^2)+U(r).
\end{equation}
利用式\eqref{ch11:eq:polar-newton}可得
$\dot E=m\vect v\cdot\vect a+U'(r)\dot r=F(r)\dot r-F(r)\dot r=0$。
对于这种没有速度相关势的模型，哈密顿量（Hamiltonian）就是上述机械能。
给定初始位置与速度，$E$ 和 $\vect L$ 就由初态确定，以后不再改变。

\subsection{束缚运动与散射}
同一组动力学方程可以回答两类不同的问题。束缚运动关心粒子是否始终停留在有限距离范围内、
是否稳定以及多久返回；散射则关心从远方入射的粒子经过力心附近后朝哪个方向离开。
图\ref{ch11:fig:bound-scatter}把这两种初始条件并列画出。

\begin{center}\begin{minipage}{\linewidth}\centering
\begin{tikzpicture}[>=Latex,scale=.84,font=\small]
 \begin{scope}[shift={(-3.8,0)}]
  \draw[curve] (-1,0) ellipse (2 and 1.732);
  \draw[curve,->,domain=45:67,samples=30] plot ({-1+2*cos(\x)},{1.732*sin(\x)});
  \fill (0,0) circle (1.8pt) node[below left=3pt] {$O$};
  \fill (1,0) circle (1.8pt) node[below right=3pt] {$P_0$};
  \draw[guide] (0,0)--(1,0) node[midway,below=5pt] {$r_0$};
  \draw[vector] (1,0)--(1,1.05) node[right=4pt] {$v_0$};
  \node at (-.7,-2.85) {(a) 有限区域内的束缚运动};
 \end{scope}
 \begin{scope}[shift={(3.3,0)},scale=.78]
  \pgfmathsetmacro{\aa}{.35}\pgfmathsetmacro{\bb}{1}
  \pgfmathsetmacro{\cc}{sqrt(\aa*\aa+\bb*\bb)}
  \pgfmathsetmacro{\cs}{\aa/\cc}\pgfmathsetmacro{\sn}{\bb/\cc}
  \draw[guide] (-3.8,1)--(3.6,1);
  \draw[curve,domain=2.0:-1.7,samples=180,variable=\u]
   plot ({\cs*(\cc-\aa*cosh(\u))-\sn*\bb*sinh(\u)},
         {\sn*(\cc-\aa*cosh(\u))+\cs*\bb*sinh(\u)});
  \fill (0,0) circle (2pt) node[below left=3pt] {$O$};
  \draw[guide] (-3.8,0)--(-2.6,0);
  \draw[<->] (-3.25,0)--(-3.25,1) node[midway,left=4pt] {$b$};
  \draw[vector] (-3.4,1)--(-2.1,1) node[midway,above=5pt] {$\vect v_\infty$};
  
  \begin{scope}[shift={(.3,-1.6)}]
   \draw[->] (0,0)--(1.35,0) node[right] {$\vect v_\infty$};
   \draw[->,DeepBlue] (0,0)--(-38.58:1.35) node[below right] {$\vect v'_\infty$};
   \draw (.58,0) arc(0:-38.58:.58);
   \node at (-19.29:.83) {$\chi$};
  \end{scope}
  \node at (-.2,-3.65) {(b) 远处速度方向的偏折};
 \end{scope}
\end{tikzpicture}
\figcap{两类有心力问题。左图的 $r_0,v_0$ 决定初始能量与角动量；右图的 $b$ 是力心到入射渐近线的垂距。}
\label{ch11:fig:bound-scatter}
\end{minipage}\end{center}

在图\ref{ch11:fig:bound-scatter}(a)中，初速度垂直于半径，初始径向速度为零；
它可能是近点或远点，却未必是圆轨道。只有径向向心条件也满足才是圆轨道。
纯粹斥力不能在无其他约束的空间维持这样的束缚轨道；有引力也不自动保证束缚或稳定。
在图(b)中，$b$ 称为冲量参数，也称瞄准距离；远处速度为 $v_\infty$ 时 $h=bv_\infty$。
散射角是远处入射、出射\emph{速度}的夹角，不是两位置矢量的夹角。
引力和斥力均可造成散射。

描述远处自由入射时，通常要求力及选定的势能随 $r\to\infty$ 衰减，并检查远场渐近运动能否定义。
“短程力”与“长程力”还涉及衰减速率及具体散射理论；平方反比引力和库仑力通常称为长程力，
但仍能研究其经典散射，不能把“必须短程”当作这里的普遍条件。
量子散射、光线在引力场中的偏折也延续了这一问题思路，不过其运动方程不能不加条件地照搬经典质点模型。

\section{有效势能与径向稳定性}
\label{ch11:sec:effective}
把 $\dot\theta=h/r^2$ 代入式\eqref{ch11:eq:energy-original}，得
\begin{equation}\label{ch11:eq:effective}
 E=\frac12m\dot r^2+U_{\rm eff}(r),\qquad
 U_{\rm eff}(r)=U(r)+\frac{mh^2}{2r^2}=U(r)+\frac{L^2}{2mr^2}.
\end{equation}
这样，二维轨道的径向变化等价于一维粒子在有效势中的运动。后一项来自横向动能，称为离心势垒；
它不是新添的实际相互作用，改变轨道角动量也会改变这条有效势曲线。
由径向方程同样得到 $m\ddot r=-U_{\rm eff}'(r)$。
允许区域满足 $E\ge U_{\rm eff}$；若等号成立且 $U_{\rm eff}'\ne0$，径向速度为零并转向。
圆轨道则必须同时满足
\begin{equation}\label{ch11:eq:circle-condition}
 U_{\rm eff}'(r_c)=0,\qquad E=U_{\rm eff}(r_c),\qquad \dot r=0.
\end{equation}

图\ref{ch11:fig:effective}取引力势 $U=-GMm/r$。此时
$r_c=h^2/(GM)$、$E_c=U_{\rm eff}(r_c)=-GMm/(2r_c)$。
令 $x=r/r_c$，便有 $U_{\rm eff}/|E_c|=x^{-2}-2x^{-1}$，纵横轴均为无量纲量。

\begin{center}\begin{minipage}{\linewidth}\centering
\begin{tikzpicture}
\begin{axis}[width=.80\linewidth,height=6.5cm,axis lines=left,
 xlabel={$r/r_c$},ylabel={$U_{\rm eff}/|E_c|$},xmin=.35,xmax=4.35,ymin=-1.45,ymax=2.0,
 xtick=\empty,ytick={-1,0,1},clip=false,
 every axis label/.style={fill=none},xlabel style={at={(axis description cs:1,-.04)},anchor=north},
 ylabel style={at={(axis description cs:-.04,1)},anchor=south}]
 \addplot[curve,domain=.42:4.25,samples=180]{1/x^2-2/x};
 \addplot[guide] coordinates {(.35,0)(4.25,0)};
 \addplot[guide] coordinates {(.35,-.5)(4.25,-.5)};
 \draw[guide] (axis cs:.585786,-1.45)--(axis cs:.585786,-.5);
 \draw[guide] (axis cs:3.414214,-1.45)--(axis cs:3.414214,-.5);
 \fill[DeepBlue] (axis cs:1,-1) circle (2pt);
 \draw[thin] (axis cs:1.1,-1.06)--(axis cs:1.65,-1.23);
 \node[anchor=west] at (axis cs:1.7,-1.25) {稳定圆轨道};
 \node[anchor=south] at (axis cs:2.8,-.47) {$E/|E_c|=-1/2$};
 \node[anchor=north] at (axis cs:.585786,-1.45) {$r_-/r_c$};
 \node[anchor=north] at (axis cs:3.414214,-1.45) {$r_+/r_c$};
\end{axis}
\end{tikzpicture}
\figcap{同一角动量下的引力有效势。水平能量线与曲线的两个交点确定近点 $r_-$ 与远点 $r_+$，最低点对应圆轨道。}
\label{ch11:fig:effective}
\end{minipage}\end{center}

图中的束缚粒子在 $r_-$ 与 $r_+$ 之间径向往返，但同时 $\dot\theta=h/r^2$ 始终为正；
所以粒子不是沿一条半径来回走。有效势先回答允许区和稳定性，精确的轨道形状仍需求 $r(\theta)$。
在圆轨道附近令 $r=r_c+\xi$，小扰动满足
\[
 m\ddot\xi=-U_{\rm eff}''(r_c)\xi+\order{\xi^2}.
\]
$U_{\rm eff}''>0$ 时是径向小振动，$U_{\rm eff}''<0$ 时扰动指数增长；
若二阶导数为零，不能立即宣布稳定，还需检查更高阶或完整方程。

\subsection{幂律力的三类有效势}
\label{ch11:sec:power}
取 $F=-kr^n$、$k>0$，这里 $n$ 是\emph{力}的指数，而不是势的指数。积分得
\[
 U(r)=\frac{k}{n+1}r^{n+1}+C\quad(n\ne-1),\qquad
 U(r)=k\ln(r/r_*)+C\quad(n=-1).
\]
参考长度 $r_*$ 只改变势能零点。$n\ne-3$ 时，由圆轨道条件有
\[
 kr_c^n=\frac{mh^2}{r_c^3},\qquad
 r_c=\left(\frac{mh^2}{k}\right)^{1/(n+3)}.
\]
对有效势再求一次导数，并用圆轨道条件消去 $k$ 或 $h$：
\begin{equation}\label{ch11:eq:power-stability}
 U_{\rm eff}''(r_c)=knr_c^{n-1}+\frac{3mh^2}{r_c^4}
 =\frac{mh^2}{r_c^4}(n+3).
\end{equation}
所以稳定圆轨道要求 $n>-3$。这不是“引力越强越稳定”：若引力在靠近力心时比 $r^{-3}$ 增长得更快，
小的向内偏离会使引力增长超过维持原角动量所需的离心项，从而继续向内偏离。

\begin{center}\begin{minipage}{\linewidth}\centering
\begin{tikzpicture}[font=\small]
 \foreach \xx/\nn/\expr in {0/-4/{.5/x^2-1/(3*x^3)},4.65/-2/{.5/x^2-1/x},9.3/1/{.5/x^2+.5*x^2}}{
  \begin{axis}[at={(\xx cm,0)},anchor=south west,width=4.6cm,height=4.7cm,
   axis lines=middle,xmin=.3,xmax=3.4,ymin=-1.1,ymax=2.2,ticks=none,
   xlabel={$r$},ylabel={$U_{\rm eff}$},clip=true,
   xlabel style={at={(axis description cs:1.02,.333333)},anchor=west},
   ylabel style={at={(axis description cs:0,1.05)},anchor=south}]
   \addplot[curve,domain=.42:3.2,samples=180]{\expr};
   \ifnum\nn=-4
    \fill (axis cs:1,.166667) circle (1.7pt);
    \node[anchor=south west] at (axis cs:1.1,.28) {$P$：不稳定};
   \else\ifnum\nn=-2
    \addplot[guide] coordinates {(.612574113,-.3)(2.72075922,-.3)};
    \fill (axis cs:1,-.5) circle (1.7pt);
    \node[anchor=north west] at (axis cs:1.1,-.60) {$P$：稳定};
   \else
    \addplot[guide] coordinates {(.618,1.5)(1.618,1.5)};
    \fill (axis cs:1,1) circle (1.7pt);
    \node[anchor=north west] at (axis cs:1.1,.87) {$P$：稳定};
   \fi\fi
  \end{axis}
 }
 \node at (2.2,-.55) {(a) $n<-3$，取 $n=-4$};
 \node at (6.85,-.55) {(b) $-3<n<-1$，取 $n=-2$};
 \node at (11.5,-.55) {(c) $n>-1$，取 $n=1$};
\end{tikzpicture}
\figcap{三类幂律引力的有效势。各面板使用各自的无量纲单位；能量线用于显示势阱内的允许区，不表示三种系统的能量相同。}
\label{ch11:fig:power-potentials}
\end{minipage}\end{center}

图\ref{ch11:fig:power-potentials}(a)中圆轨道处是势垒顶而非势阱底；
向两侧扰动可导致向力心坠落或向远处逃离。图(b)的有效势在无穷远趋于零，
负能量的适当范围可束缚，非负能量则可通向无穷远。图(c)的势能在远处无限增大，
有限能量运动的外侧也受到限制。$n=-1$ 的对数势同样有内侧离心势垒和远处上升的势能，
应直接用对数式讨论，不能代入含 $1/(n+1)$ 的公式。

当 $n=-3$ 时，
\[
 U_{\rm eff}=\frac{mh^2-k}{2r^2}+C.
\]
若 $mh^2=k$，有效势为常数，各个半径都可在 $\dot r=0$ 时维持圆周运动，
但没有径向回复力；加入微小径向速度便会漂移，不是稳定势阱中的振动。
若 $mh^2\ne k$，有效势单调，不存在圆轨道。这个临界情形应与一般幂律分开。

\section{从径向运动求时间与轨道}
\label{ch11:sec:quadrature}
将有效势能量式解出径向速度，必须保留方向分支：
\[
 \dot r=\pm\sqrt{\frac2m[E-U_{\rm eff}(r)]}.
\]
出向段取正号、入向段取负号，转向点处连接两支。
从 $r_a$ 到 $r_b$ 的一段运动，其时间可按沿途方向取
\[
 \Delta t=\int_{r_a}^{r_b}\frac{\dd r}{\dot r}.
\]
对于图\ref{ch11:fig:effective}中两个简单转向点之间的束缚运动，径向周期为
\begin{equation}\label{ch11:eq:radial-period}
 T_r=2\int_{r_-}^{r_+}\frac{\dd r}{\sqrt{2[E-U_{\rm eff}(r)]/m}}.
\end{equation}
因从近点到远点只占一个径向周期的一半，前面有因子2。
端点虽使分母为零，但简单转向点附近积分形如 $\int\dd r/\sqrt{|r-r_\pm|}$，通常收敛。
恰好圆轨道的两个端点重合，应改用小振动极限，而不是直接把积分上下限设成同一点。

用 $\dot\theta=h/r^2$ 除以 $\dot r$，得到
\begin{equation}\label{ch11:eq:angle-integral}
 \theta_b-\theta_a=\pm\int_{r_a}^{r_b}
 \frac{h\,\dd r}{r^2\sqrt{2[E-U_{\rm eff}(r)]/m}}.
\end{equation}
这里正负号也由出向/入向段决定。一般位置之间的积分端点是这两个位置，
\emph{不是必须取转向点}。只有求近点到远点夹角、完整径向周期或散射全过程时，才按问题选择转向点及无穷远。
例如每个径向周期内转过的方位角为
\[
 \Delta\theta_r=2\int_{r_-}^{r_+}
 \frac{h\,\dd r}{r^2\sqrt{2[E-U_{\rm eff}(r)]/m}}.
\]
若 $\Delta\theta_r/(2\pi)$ 是有理数，径向状态和方位角可以在整数个周期后同时重复；
若为无理数，轨道一般不闭合。径向周期和绕力心转一圈所需时间不能先验地当成同一个量。

\subsection{散射角与冲量参数}
远方取 $U(\infty)=0$，则 $E=mv_\infty^2/2$、$h=bv_\infty$。
设散射轨道有一个最近距离 $r_m$，由 $E=U_{\rm eff}(r_m)$ 求出。
从最近点到无穷远的方位角变化为
\[
 \Phi=\int_{r_m}^{\infty}
 \frac{b\,\dd r}{r^2\sqrt{1-b^2/r^2-2U(r)/(mv_\infty^2)}}.
\]
不受力的直线轨道有 $\Phi=\pi/2$。对于不绕转多圈的通常单次散射，
远处速度的偏折角为 $\chi=|2\Phi-\pi|$；引力向力心弯折、斥力背离力心弯折。
若轨道绕力心多圈或有多个转向结构，需保留累积方位角，并按实际入出射方向取夹角，不能丢掉绕转次数。

\subsection{比内方程的换元}
轨道形状用 $r(\theta)$ 描述比先求 $r(t),\theta(t)$ 再消去时间更直接。
令 $u(\theta)=1/r$，以下撇号表示对 $\theta$ 求导。首先
\[
 \dot\theta=hu^2,\qquad
 \frac{\dd r}{\dd\theta}=-\frac{u'}{u^2},\qquad
 \frac{\dd^2r}{\dd\theta^2}=-\frac{u''}{u^2}+\frac{2(u')^2}{u^3}.
\]
由链式法则分别得到
\[
 \dot r=-\frac{u'}{u^2}\,hu^2=-hu',\qquad
 \ddot r=-hu''\dot\theta=-h^2u^2u'',\qquad
 r\dot\theta^2=h^2u^3.
\]
代入径向牛顿方程，整理为比内方程（Binet equation）：
\begin{equation}\label{ch11:eq:binet}
 u''+u=-\frac{F(1/u)}{mh^2u^2}
       =-\frac{mF(1/u)}{L^2u^2}.
\end{equation}
左边和右边的量纲均为长度的倒数。这一换元把时间消去，但不改变轨道的物理条件；
径向退化运动 $h=0$ 仍应直接用时间方程处理。

\subsection{积分上下限和运动分支如何选择}
上面的积分式解决三个不同问题：时间积分用来问“多久到达”，角度积分用来问“到达何处”，比内方程用来直接求轨道形状。列式之前，先由 $E\ge U_{\rm eff}(r)$ 找允许区间，再用初始径向速度选入向或出向分支。不能因根号约定为正，就把正在靠近力心的 $\dot r$ 也取正。

\Needspace{16\baselineskip}\begin{origexamplebox}{从静止径向落向引力中心的时间}\label{ch11:ex:teaching-radial-fall}
质量可忽略的质点从距固定点质量 $M$ 为 $r_0$ 处静止释放，仅受万有引力。求降到 $r_1$ 的时间，其中 $0<r_1<r_0$，并讨论形式上的 $r_1\to0$ 极限。把引力源看作点质量。

\ThQRadialFall
\end{origexamplebox}
\Needspace{6\baselineskip}\begin{xesolution}
\ThARadialFall

初速度为零，故角动量为零，运动是纯径向的，不可使用含 $1/L$ 的轨道公式。记 $\mu=GM$，单位质量机械能守恒给出
\[
 \tfrac12\dot r^2-\frac\mu r=-\frac\mu{r_0}.
\]
质点向内落，选负根 $\dot r=-\sqrt{2\mu(1/r-1/r_0)}$。为使时间积分为正，写成
\[
 t=\int_{r_1}^{r_0}\frac{\dd r}{\sqrt{2\mu(1/r-1/r_0)}}.
\]
令 $r=r_0\cos^2u$，$u_1=\arccos\sqrt{r_1/r_0}$。则 $|\dd r|=2r_0\sin u\cos u\,\dd u$，分母为 $\sqrt{2\mu/r_0}\tan u$，故
\[
 t=\sqrt{\frac{2r_0^3}{\mu}}\int_0^{u_1}\cos^2u\,\dd u
 =\sqrt{\frac{r_0^3}{2\mu}}\left(u_1+\sin u_1\cos u_1\right).
\]
回到半径变量，括号中的第二项为 $\sqrt{(r_1/r_0)(1-r_1/r_0)}$。当 $r_1\to r_0$，时间趋零；当 $r_1\to0$，形式极限为 $\pi\sqrt{r_0^3/(8GM)}$，量纲为时间。
现实中有半径的天体会先与质点接触；点质量模型在中心还有奇点，因此有限的形式时间不意味着质点能无条件穿过中心继续沿同一公式运动。
\end{xesolution}

\section{平方反比力与万有引力二体问题}
\label{ch11:sec:kepler}
经典平方反比力问题是天体力学的重要起点。统一写成 $F=-m\kappa/r^2$，
则引力的 $\kappa=GM>0$；库仑相互作用可取
$\kappa=-Qq/(4\pi\varepsilon_0m)$，同号电荷给 $\kappa<0$。
有些文献用 $\sqrt{GM}$ 作为引力强度参数；本章始终直接使用 $GM$，避免平方与符号混淆。
对 $\kappa>0$，式\eqref{ch11:eq:binet}成为常系数方程
\[
 u''+u=\frac{\kappa}{h^2}.
\]
特解为常数，齐次解为正弦与余弦的组合。将两积分常数合并为振幅和初相位：
\begin{equation}\label{ch11:eq:conic}
 u=\frac{\kappa}{h^2}[1+e\cos(\theta-\theta_0)],\qquad
 r=\frac{p}{1+e\cos(\theta-\theta_0)},\quad p=\frac{h^2}{\kappa}.
\end{equation}
$e\ge0$ 是偏心率，$p$ 为半通径，$\theta_0$ 取近点方向；$e=0$ 时是圆，近点方向不唯一。
物理轨道只取分母为正的区域。$0<e<1$ 为椭圆，$e=1$ 为抛物线，$e>1$ 为双曲线。
这称为开普勒问题。对斥力 $\kappa<0$，更方便把解写成
$u=-|\kappa|/h^2+C\cos(\theta-\theta_0)$，取 $u>0$ 的双曲线支；
不能仍把 $p=h^2/\kappa$ 当成正长度并沿用引力的全部范围。

\subsection{有效势积分给出能量与偏心率}
比内方程容易给出圆锥曲线，却还没有把 $e$ 与初始能量联系起来。
令比机械能 $\varepsilon=E/m$，将 $\dot r=-hu'$ 代入引力能量式：
\[
 \varepsilon=\tfrac12h^2(u')^2+\tfrac12h^2u^2-\kappa u.
\]
配方而不是省略交叉项，得
\[
 (u')^2+\left(u-\frac{\kappa}{h^2}\right)^2
 =\frac{\kappa^2}{h^4}+\frac{2\varepsilon}{h^2}.
\]
记右边平方根为 $C$，则
\[
 \frac{\dd(u-\kappa/h^2)}{\sqrt{C^2-(u-\kappa/h^2)^2}}=\pm\dd\theta.
\]
积分并把常相位吸收进 $\theta_0$，恢复式\eqref{ch11:eq:conic}，同时得到
\begin{equation}\label{ch11:eq:energy-eccentricity}
 e^2=1+\frac{2\varepsilon h^2}{\kappa^2}
     =1+\frac{2EL^2}{m^3\kappa^2}.
\end{equation}
因此负、零、正能量分别对应椭圆、抛物线、双曲线；这依赖引力势零点取在无穷远。
对于椭圆，近远点之和为 $2a$，
\[
 r_p=\frac p{1+e}=a(1-e),\qquad
 r_a=\frac p{1-e}=a(1+e),\qquad p=a(1-e^2).
\]
代回能量关系得
\begin{equation}\label{ch11:eq:vis-viva}
 E=-\frac{m\kappa}{2a},\qquad
 v^2=\kappa\left(\frac2r-\frac1a\right).
\end{equation}
后一式称为活力公式。对于双曲线若用正的几何半实轴 $a_h=p/(e^2-1)$，
则 $E=+m\kappa/(2a_h)$；抛物线用 $E=0$ 单独处理。

于是给定初始 $r_0,v_{r0},v_{\theta0}$ 后，可以依次计算
\[
 h=r_0v_{\theta0},\qquad
 \varepsilon=\tfrac12(v_{r0}^2+v_{\theta0}^2)-\frac\kappa{r_0},\qquad
 p=\frac{h^2}{\kappa},\quad e=\sqrt{1+\frac{2\varepsilon h^2}{\kappa^2}}.
\]
当 $e>0$ 时，由
$e\cos(\theta-\theta_0)=p/r_0-1$ 和
$e\sin(\theta-\theta_0)=hv_{r0}/\kappa$ 共同确定初相位象限，
不能只用反余弦丢掉入向与出向之别。

\subsection{两颗星都运动时怎样使用上述公式}
设两颗星质量为 $m_1,m_2$，没有外力，定义
\[
 M=m_1+m_2,\quad
 \vect R=\frac{m_1\vect r_1+m_2\vect r_2}{M},\quad
 \vect r=\vect r_1-\vect r_2,\quad \mu=\frac{m_1m_2}{M}.
\]
$\mu$ 称为约化质量。相加两个牛顿方程得 $M\ddot{\vect R}=0$；相减则得
\[
 \ddot{\vect r}=-GM\frac{\vect r}{r^3},\qquad
 \mu\ddot{\vect r}=-Gm_1m_2\frac{\vect r}{r^3}.
\]
因此前述固定中心方程仍适用，但运动的是\emph{两星的相对位置}，引力参数为 $GM$。
由 $\vect r_1=\vect R+(m_2/M)\vect r$、$\vect r_2=\vect R-(m_1/M)\vect r$，展开并相消交叉项：
\[
 K=\tfrac12M|\dot{\vect R}|^2+\tfrac12\mu|\dot{\vect r}|^2,\qquad
 E_{\rm rel}=\tfrac12\mu|\dot{\vect r}|^2-\frac{Gm_1m_2}{r}.
\]
相对角动量为 $\vect L_{\rm rel}=\mu\vect r\times\dot{\vect r}$，故 $h=L_{\rm rel}/\mu$。
在质心系，两颗星的轨道与相对轨道相似、比例分别为 $m_2/M$ 与 $m_1/M$，且位于质心两侧。
只有 $m_2\ll m_1$ 时才能近似把 $GM$ 换为较大星体的 $Gm_1$。

\subsection{圆轨道、逃逸速度与周期}
圆轨道由 $mv_c^2/r=GMm/r^2$ 得
\[
 v_c=\sqrt{GM/r},\qquad T_c=\frac{2\pi r}{v_c}=2\pi\sqrt{r^3/(GM)}.
\]
若从距离 $r$ 处恰好逃到无穷远且末速为零，则初态总能量为零：
$mv_e^2/2-GMm/r=0$，所以 $v_e=\sqrt{2GM/r}=\sqrt2v_c$。
对椭圆不能直接用周长除以某处速度，因为速度不恒定。设短半轴为 $b_e=a\sqrt{1-e^2}$，
用式\eqref{ch11:eq:area-law}的面积速度除去椭圆面积：
\[
 T=\frac{\pi ab_e}{h/2}=\frac{2\pi ab_e}{\sqrt{GMp}}
 =2\pi\sqrt{\frac{a^3}{GM}}.
\]
故开普勒第三定律为
\begin{equation}\label{ch11:eq:third-law}
 \boxed{T^2=\frac{4\pi^2a^3}{G(m_1+m_2)}}.
\end{equation}
其中 $a$ 是相对椭圆的半长轴；若给某一颗星相对质心的半长轴，必须先按质量比换算。

\Needspace{16\baselineskip}\begin{origexamplebox}{近地圆轨道和逃逸速度}\label{ch11:ex:01}
忽略大气和地球自转，地球半径为 $R$、地表重力加速度为 $g$。求近地圆轨道速度、周期和第二宇宙速度。

\ThQOrbit
\end{origexamplebox}
\Needspace{6\baselineskip}\begin{xesolution}
\ThAOrbit

地表重力满足 $mg=GMm/R^2$，消去试探物体质量得 $GM=gR^2$。
近地圆轨道径向方程为 $mv_c^2/R=mg$，因此
\[
 \boxed{v_c=\sqrt{gR}},\qquad
 \boxed{T=\frac{2\pi R}{v_c}=2\pi\sqrt{\frac Rg}}.
\]
恰好逃逸是能量阈值而非“重力突然消失”：途中仍受引力，动能逐渐减小。
令初态机械能等于无穷远静止时的零能量，
\[
 \frac12mv_e^2-\frac{GMm}{R}=0
 \quad\Longrightarrow\quad\boxed{v_e=\sqrt{2gR}}.
\]
$gR$ 的量纲为速度平方、$R/g$ 为时间平方。取 $R=6.37\times10^6\,\mathrm m$、
$g=9.81\,\mathrm{m/s^2}$，分别约为 $7.9\,\mathrm{km/s}$、$84.4\,\mathrm{min}$ 和 $11.2\,\mathrm{km/s}$。
这是理想近地极限；实际卫星不能贴地飞行，需避开地形与大气阻力。
\end{xesolution}

\Needspace{16\baselineskip}\begin{origexamplebox}{椭圆轨道远近点速度}\label{ch11:ex:02}
卫星绕质量为 $M$ 的天体作椭圆运动，近点和远点距离分别为 $r_p,r_a$。求两点速度之比，并给出速度本身。

\ThQApsides
\end{origexamplebox}
\Needspace{6\baselineskip}\begin{xesolution}
图\ref{ch11:fig:apsis}中力心在椭圆的焦点而非中心。两拱点有 $\dot r=0$，速度均垂直于位置矢量。
\begin{center}\begin{minipage}{\linewidth}\centering
\begin{tikzpicture}[>=Latex,scale=.78,font=\small]
 \draw[curve] (0,0) ellipse (3 and 1.8);
 \coordinate (S) at (2.4,0);
 \fill (S) circle (2pt) node[above left=4pt] {$S$};
 \fill (3,0) circle (2pt) node[below right=4pt] {$P$};
 \fill (-3,0) circle (2pt) node[below left=4pt] {$A$};
 \draw[vector] (3,0)--(3,1.3) node[right=4pt] {$v_p$};
 \draw[vector] (-3,0)--(-3,-.7) node[left=4pt] {$v_a$};
 \draw[guide] (-3,0)--(3,0);
 \draw[<->] (2.4,-2.2)--(3,-2.2) node[midway,below=5pt] {$r_p$};
 \draw[<->] (-3,-2.2)--(2.4,-2.2) node[midway,below=5pt] {$r_a$};
\end{tikzpicture}
\figcap{近点 $P$ 与远点 $A$ 的速度沿同一绕行方向。两距离之和等于椭圆长轴 $2a$。}\label{ch11:fig:apsis}
\end{minipage}\end{center}
角动量守恒给出 $r_pv_p=r_av_a$，故 $v_p/v_a=r_a/r_p$。
这还没有给出速度的绝对大小。利用两点机械能相等，
\[
 \tfrac12v_p^2-\frac{GM}{r_p}
 =\tfrac12v_a^2-\frac{GM}{r_a}.
\]
代入 $v_a=(r_p/r_a)v_p$，在非圆情形消去 $r_a-r_p$，
\[
 v_p^2\frac{(r_a-r_p)(r_a+r_p)}{2r_a^2}
 =GM\frac{r_a-r_p}{r_pr_a}.
\]
于是
\[
 \boxed{v_p=\sqrt{\frac{2GMr_a}{r_p(r_p+r_a)}}},\qquad
 \boxed{v_a=\sqrt{\frac{2GMr_p}{r_a(r_p+r_a)}}},\qquad
 \boxed{\frac{v_p}{v_a}=\frac{r_a}{r_p}}.
\]
也可把图中 $a=(r_p+r_a)/2$ 代入活力公式\eqref{ch11:eq:vis-viva}复核。
圆轨道 $r_p=r_a=r$ 虽不能在上一步除以零，但由连续极限或向心方程都得 $v=\sqrt{GM/r}$。
\end{xesolution}

\section{轨道能量、变速与拱线进动}
椭圆总机械能 $E=-GMm/(2a)$ 只依赖半长轴，角动量则还区分形状。
在某点瞬时变速时，位置和势能未变，能量改变量来自动能，角动量改变量来自速度的切向分量。
所以只给“增加多少能量”通常不足以确定一条新轨道，还要给变速方向。

从半径 $r_0$ 的\emph{圆轨道}出发，若只作沿原切向的瞬时变速且不反向，则 $\dot r=0$ 仍成立。
新速度 $v$ 满足
\[
 \ddot r=\frac{v^2}{r_0}-\frac{GM}{r_0^2}.
\]
若 $v_c<v<v_e$，该点是新椭圆近点；若 $0<v<v_c$，该点是远点；$v=v_c$ 保持圆轨道。
$v=v_e$ 是抛物轨道、$v>v_e$ 是双曲轨道，都不再是束缚椭圆。
若 $v=0$，则退化为径向坠落。若从一般椭圆上 $\dot r\ne0$ 的位置变速，
仅改变切向速度不会消去径向速度，该位置不能因此成为近点或远点。

\subsection{一次变速先确定新轨道，再讨论怎样接续}
瞬时变速只改变速度，不改变该瞬间的位置。此后若发动机关闭，卫星沿由新能量、新角动量确定的轨道运动；不能把“速度变大”直接翻译成“立即进入更大的圆”。较高圆轨道的圆周速度虽然更低，从较低圆轨道转移过去却通常要先顺行加速，因为首先要提高轨道能量。

\Needspace{16\baselineskip}\begin{origexamplebox}{两个共面圆轨道之间的双脉冲转移}\label{ch11:ex:teaching-transfer}
卫星原在半径 $r_1$ 的圆轨道上，拟转入同平面、同绕行方向、半径 $r_2>r_1$ 的圆轨道。采用一条近点为 $r_1$、远点为 $r_2$ 的椭圆转移轨道，在两端各施加一次沿切向的瞬时速度增量。忽略卫星质量及推进过程的有限持续时间，求两次增量与转移时间。

\ThQTransfer
\end{origexamplebox}
\Needspace{6\baselineskip}\begin{xesolution}
\ThATransfer

图~\ref{ch11:fig:apsis} 的 $P,A$ 分别对应转移椭圆的近点、远点，故半长轴 $a=(r_1+r_2)/2$。记 $\mu=GM$。转移过程中单位质量能量为 $-\mu/(2a)$，在任一点有
\[
 \frac{v^2}{2}-\frac\mu r=-\frac\mu{2a},
 \qquad v^2=\mu\left(\frac2r-\frac1a\right).
\]
这先由能量守恒给出端点速度，尚未把它们误认为圆周速度。在 $P$ 点：
\[
 v_P=\sqrt{\frac{2\mu r_2}{r_1(r_1+r_2)}},\qquad
 \Delta v_1=v_P-\sqrt{\frac\mu{r_1}}>0.
\]
第一次顺行加速，把圆轨道切换为转移椭圆。飞到 $A$ 点时：
\[
 v_A=\sqrt{\frac{2\mu r_1}{r_2(r_1+r_2)}},\qquad
 \Delta v_2=\sqrt{\frac\mu{r_2}}-v_A>0.
\]
第二次仍是顺行加速，因为转移椭圆远点速度比当地圆周速度更小。两端满足 $r_1v_P=r_2v_A$，核验了无推力飞行段的角动量守恒；两次脉冲处卫星自身角动量则改变，不能跨过脉冲继续当作守恒。
从近点到远点用半个椭圆周期，因此
\[
 t_{\rm tr}=\pi\sqrt{\frac{a^3}{\mu}}
 =\pi\sqrt{\frac{(r_1+r_2)^3}{8\mu}}.
\]
$r_2\to r_1$ 时两次速度增量均趋零；形式上的飞行时间趋于原圆轨道半周期，表示仍按“飞到对面”的约定计时，而不是说留在原轨道必须等待半周期。结果的速度与时间量纲分别由 $\sqrt{\mu/r}$、$\sqrt{r^3/\mu}$ 给出。
\end{xesolution}

\begingroup\tcbset{breakable/.default=true}\begin{starread}{近圆轨道的径向频率与拱线}
\label{ch11:sec:precession}
设圆轨道平均角频率为 $\Omega=h/r_c^2$，径向小振动频率
$\kappa_r=\sqrt{U_{\rm eff}''(r_c)/m}$。每个径向周期转过
$\Delta\theta_r\simeq2\pi\Omega/\kappa_r$。
对于力指数为 $n$ 的幂律引力，式\eqref{ch11:eq:power-stability}给
$\kappa_r^2=(n+3)\Omega^2$。
平方反比引力 $n=-2$ 有 $\kappa_r=\Omega$，相邻近日点方向相同；
各向同性谐振子 $n=1$ 则有 $\kappa_r=2\Omega$，相邻最近点相差 $\pi$，
是同一条固定椭圆短轴上的另一端，并不是长轴持续进动。
因此“两个频率不相等”不能直接推出拱线转动，必须检查方位角及方向的等价关系。

在接近开普勒势的弱扰动中，$\kappa_r$ 接近 $\Omega$，每径向周期超过 $2\pi$ 的小角
$\delta\varpi=2\pi(\Omega/\kappa_r-1)$ 对应平均进动速率
$\dot\varpi=\delta\varpi/(2\pi/\kappa_r)=\Omega-\kappa_r$。
这个公式的使用背景是近圆、近开普勒轨道，不能不加解释地套到所有幂律势。
在通常光滑有心势的条件下，使所有束缚轨道都闭合的平方反比引力和各向同性谐振子是两个重要特例。
\end{starread}\endgroup

\Needspace{22\baselineskip}
\section{天体运动综合问题}

\Needspace{16\baselineskip}\begin{origexamplebox}{哈雷彗星的轨道状态}
\label{ch11:ex:03}
哈雷彗星绕太阳沿椭圆轨道逆时针运动，周期为 $T=76.1\,\mathrm y$。1986年经过近日点
$P_0$ 时，日彗距离为 $r_0=0.590\,\mathrm{AU}$；此后到达 $P$ 点，位置矢量
$\overrightarrow{SP}$ 与近日点方向 $\overrightarrow{SP_0}$ 的夹角为 $\theta=72^\circ$。
取 $1\,\mathrm{AU}=1.5\times10^{11}\,\mathrm m$。求此时的日彗距离、速度大小和速度方向
（以速度矢量与射线 $SP_0$ 的逆时针夹角表示）。


\begin{center}\begin{minipage}{\linewidth}\centering
\begin{tikzpicture}[scale=.90,>=Latex]
  \coordinate (S) at (0,0);
  \draw[axis,->] (0,0)--(4.3,0) node[right] {近日点方向};
  \draw[curve,samples=180,smooth,variable=\a,domain=-104:104]
    plot ({2.0*1.16062*cos(\a)/(1+.96715*cos(\a))},
          {2.0*1.16062*sin(\a)/(1+.96715*cos(\a))});
  \fill (S) circle (2.6pt) node[below left] {太阳 $S$};
  \coordinate (P0) at (1.18,0);
  \fill[DeepBlue] (P0) circle (2.2pt) node[below right] {$P_0$};
  \coordinate (P) at ({1.78712*cos(72)},{1.78712*sin(72)});
  \draw[guide] (S)--(P);
  \fill[DeepBlue] (P) circle (2.4pt) node[above right] {$P$};
  \draw[->,TextGray] (.70,0) arc[start angle=0,end angle=72,radius=.70];
  \node at (36:.95) {$\theta$};
\end{tikzpicture}
\figcap{题设轨道的近日点附近示意。$S$ 为太阳，$P_0$ 为近日点，$\theta=72^\circ$。}
\label{ch11:fig:halley}
\end{minipage}\end{center}
\end{origexamplebox}
\Needspace{6\baselineskip}\begin{xesolution}
图\ref{ch11:fig:halley-analysis}中 $SP_0$ 是位置角的零方向，速度沿轨道切线。
\begin{center}\begin{minipage}{\linewidth}\centering
\begin{tikzpicture}[scale=.90,>=Latex]
  \coordinate (S) at (0,0);
  \draw[axis,->] (0,0)--(4.3,0) node[right] {近日点方向};
  \draw[curve,samples=180,smooth,variable=\a,domain=-104:104]
    plot ({2.0*1.16062*cos(\a)/(1+.96715*cos(\a))},
          {2.0*1.16062*sin(\a)/(1+.96715*cos(\a))});
  \fill (S) circle (2.6pt) node[below left] {太阳 $S$};
  \coordinate (P0) at (1.18,0);
  \fill[DeepBlue] (P0) circle (2.2pt) node[below right] {$P_0$};
  \coordinate (P) at ({1.78712*cos(72)},{1.78712*sin(72)});
  \draw[guide] (S)--(P);
  \fill[DeepBlue] (P) circle (2.4pt) node[above right] {$P$};
  \draw[->,DeepBlue,very thick] (P)--++(126.70:1.25) node[above left] {$\vect v$};
  \draw[->,TextGray] (.70,0) arc[start angle=0,end angle=72,radius=.70];
  \node at (36:.95) {$\theta$};
\end{tikzpicture}
\figcap{近日点方向取作极轴。图中 $\theta$ 是位置角，而速度方向还要由径向、切向两个分量共同确定。}
\label{ch11:fig:halley-analysis}
\end{minipage}\end{center}
以地球轨道为单位使用开普勒第三定律，太阳质量占主导时有
\[
 \left(\frac{T}{1\,\mathrm y}\right)^2
 =\left(\frac{a}{1\,\mathrm{AU}}\right)^3,
 \qquad
 a=(76.1)^{2/3}\,\mathrm{AU}=17.9579\,\mathrm{AU}.
\]
近日点满足 $r_0=a(1-e)$，所以
\[
 e=1-\frac{r_0}{a}=0.967145,
 \qquad
 p=a(1-e^2)=1.16062\,\mathrm{AU}.
\]
图中 $SP_0$ 已选作真近点角的零方向，故在 $\theta=72^\circ$ 处
\[
 \boxed{r=\frac{p}{1+e\cos\theta}=0.893562\,\mathrm{AU}}
 \approx1.34\times10^{11}\,\mathrm m.
\]

在 $\mathrm{AU}$、年为单位时，太阳引力参数为
$\mu_\odot=GM_\odot=4\pi^2\,\mathrm{AU^3/y^2}$。由活力公式
\[
 v=\sqrt{\mu_\odot\left(\frac2r-\frac1a\right)}
   =9.28243\,\mathrm{AU/y},
\]
取一年为 $365.25$ 天，把 $\mathrm{AU/y}$ 换算成 $\mathrm{km/s}$，即
\[
 \boxed{v=44.0\,\mathrm{km/s}}.
\]
只给速度大小还不能确定图中的箭头。令单位质量角动量
$h=r^2\dot\theta=\sqrt{\mu_\odot p}$。对轨道方程求导，
\[
 \frac{\dd r}{\dd\theta}=\frac{pe\sin\theta}{(1+e\cos\theta)^2},\qquad
 \dot r=\frac{\dd r}{\dd\theta}\frac h{r^2}=\frac h p e\sin\theta,\qquad
 r\dot\theta=\frac h r.
\]
利用 $h/p=\mu_\odot/h$，沿径向外和逆时针切向的分量分别为
\[
 v_r=\frac{\mu_\odot}{h}e\sin\theta=5.36456\,\mathrm{AU/y},
 \qquad
 v_\theta=\frac{\mu_\odot}{h}(1+e\cos\theta)=7.57529\,\mathrm{AU/y}.
\]
因 $v_r>0$，彗星正在远离太阳。速度相对径向 $SP$ 再逆时针偏转
$\arctan(v_\theta/v_r)=54.695^\circ$，所以相对 $SP_0$ 的方向角为
\[
 \boxed{\psi=72^\circ+54.695^\circ=126.695^\circ}.
\]
代回 $v^2=v_r^2+v_\theta^2$ 可恢复前述活力公式数值，方向与大小相互校验。
\end{xesolution}

\Needspace{16\baselineskip}\begin{origexamplebox}{双星系统中的瞬时对称失质量}
\label{ch11:ex:04}
质量分别为 $M$ 和 $m$（$M>m$）的两颗星相距 $d$，在相互引力作用下绕共同质心作圆周运动。
某次超新星爆发使质量较大的星体瞬时、各向同性地失去质量 $\Delta M$；抛出物不碰撞另一颗星，
且爆发不对残余星体施加反冲。求爆发后两颗残余星体仍保持束缚的条件。

\ThQMassLoss
\end{origexamplebox}
\Needspace{6\baselineskip}\begin{xesolution}
\ThAMassLoss

爆发前相对运动的圆轨道速度满足
\[
 v_{\rm rel}^2=\frac{G(M+m)}{d}.
\]
爆发是瞬时的，位置与两颗残余星体的速度均来不及改变；爆发后的引力质量和却变为
$M-\Delta M+m$。于是单位约化质量的相对轨道能量为
\begin{align*}
 \varepsilon'
 &=\frac12v_{\rm rel}^2-\frac{G(M-\Delta M+m)}{d}\\
 &=\frac{G}{2d}\,[2\Delta M-(M+m)].
\end{align*}
实际失质量还应满足 $0\le\Delta M<M$。残余双星保持束缚等价于 $\varepsilon'<0$，故
\[
 \boxed{\Delta M<\frac{M+m}{2}}.
\]
等号对应抛物线临界逃逸；超过该值为双曲线分离。这个判据依赖“瞬时、对称、无反冲”，若爆发给残余星体一个踢速度，
还必须把该速度矢量加入爆发瞬间的相对速度后重新计算能量。
\end{xesolution}

\Needspace{16\baselineskip}\begin{origexamplebox}{拉普拉斯--龙格--楞次矢量}
\label{ch11:ex:05}
质量为 $m$ 的质点在质量为 $M$ 的固定中心产生的平方反比引力中运动。证明
拉普拉斯--龙格--楞次矢量（Laplace--Runge--Lenz vector）
\[
 \vect e=\frac{\vect p\times\vect L}{GMm^2}-\frac{\vect r}{r}
\]
守恒，并由此推出轨道的圆锥曲线方程
\[
 r=\frac{p_0}{1+e\cos\theta}.
\]

\ThQLRL
\end{origexamplebox}
\Needspace{6\baselineskip}\begin{xesolution}
\ThALRL

这里 $\vect p=m\vect v$、$\vect L=\vect r\times\vect p$，并记
$\uvec r=\vect r/r$。引力为
\[
 \vect F=-\frac{GMm}{r^3}\vect r,
\]
所以 $\dot{\vect L}=\vect r\times\vect F=0$。另一方面，利用三重积恒等式，
\begin{align*}
 \frac{\dd}{\dd t}(\vect p\times\vect L)
 &=\vect F\times\vect L\\
 &=-\frac{GMm}{r^3}\,
   \vect r\times(\vect r\times\vect p)\\
 &=GMm\left(\frac{\vect p}{r}
 -\frac{\vect r(\vect r\cdot\vect p)}{r^3}\right)
 =GMm^2\frac{\dd\uvec r}{\dd t}.
\end{align*}
因此 $\dot{\vect e}=0$，所给矢量确为常矢量。

取极轴沿 $\vect e$，令 $\theta$ 为 $\vect e$ 与 $\vect r$ 的夹角。将守恒式点乘 $\uvec r$，并注意
\[
 (\vect p\times\vect L)\cdot\uvec r
 =\frac1r\vect L\cdot(\vect r\times\vect p)=\frac{L^2}{r},
\]
便有
\[
 e\cos\theta=\frac{L^2}{GMm^2r}-1.
\]
整理得
\[
 \boxed{r=\frac{p_0}{1+e\cos\theta}},
 \qquad
 \boxed{p_0=\frac{L^2}{GMm^2}}.
\]
这里取 $L\ne0$。$0<e<1,e=1,e>1$ 分别给出椭圆、抛物线、双曲线；
矢量的模就是轨道偏心率，方向指向近心点。$e=0$ 时轨道为圆，$\vect e=0$，
不能沿零矢量选极轴，任取平面内固定方向即可。$L=0$ 的纯径向运动直接用径向能量方程，
不将 $p_0=0$ 代入非退化轨道公式。$\vect e$ 无量纲，$p_0$ 具有长度量纲。
\end{xesolution}

\Needspace{16\baselineskip}\begin{origexamplebox}{星际尘埃造成的引力阻力}
\label{ch11:ex:06}
宇宙飞船质量为 $M_s$、几何半径为 $r_s$，以速度 $v$ 穿过密度为 $\rho$ 的均匀稀薄尘埃。
尘埃微粒质量远小于飞船质量，远处可视为平行来流；只有冲量参数不大于 $R_d$ 的尘埃计入环境范围。

（1）先把飞船视为固定的引力中心，求尘埃被散射到 $\theta$ 至 $\theta+\dd\theta$ 时的微分散射截面
$\dd\sigma/\dd\Omega$；（2）假设直接撞击的尘埃被飞船吸收，并且引力散射均为小角度，求直接碰撞阻力与
掠过飞船的尘埃造成的引力阻力。

\ThQDust
\end{origexamplebox}
\Needspace{6\baselineskip}\begin{xesolution}
设尘埃的冲量参数为 $b$，即远处入射直线与引力中心的距离，而不是最近距离。
远处单位质量能量为 $v^2/2$，单位质量角动量为 $bv$。代入偏心率公式得到
\[
 e=\sqrt{1+\frac{b^2v^4}{G^2M_s^2}}=\sqrt{1+\frac{b^2}{a_g^2}},
 \qquad a_g=\frac{GM_s}{v^2}.
\]
在无穷远处 $1+e\cos\nu_\infty=0$；入射与出射速度的夹角为
$\theta=2\nu_\infty-\pi$，故 $\sin(\theta/2)=1/e$，从而
\[
 \tan\frac\theta2=\frac{a_g}{b},
 \qquad
 a_g=\frac{GM_s}{v^2}.
\]
图\ref{ch11:fig:scattering}左侧实线是双曲线轨迹，右侧单独比较远处的速度矢量，
不把近处尚在变化的切线方向当作最终散射方向。
\begin{center}\begin{minipage}{\linewidth}\centering
\begin{tikzpicture}[scale=.87,>=Latex]
  \pgfmathsetmacro{\ag}{.35}\pgfmathsetmacro{\cc}{sqrt(1+\ag*\ag)}
  \pgfmathsetmacro{\cs}{\ag/\cc}\pgfmathsetmacro{\sn}{1/\cc}
  \pgfmathsetmacro{\chi}{2*atan(\ag)}
  \draw[guide] (-4.1,0)--(3.4,0);
  \draw[guide] (-4.1,1)--(-.2,1);
  \draw[<->] (-3.9,0)--(-3.9,1) node[midway,left] {$b$};
  \draw[curve,domain=2.05:-1.8,samples=180,variable=\u,->]
   plot ({\cs*(\cc-\ag*cosh(\u))-\sn*sinh(\u)},
         {\sn*(\cc-\ag*cosh(\u))+\cs*sinh(\u)});
  \filldraw[draw=DeepBlue,fill=DeepBlue!12] (0,0) circle(.14);
  \node[below] at (0,-.18) {$M_s$};
  \draw[->,DeepBlue] (-3.4,1.4)--(-2.2,1.4) node[midway,above] {远处入射};
  \node at (-.3,-1.65) {轨迹：$a_g/b=0.35$};
  \begin{scope}[shift={(5,1)}]
   \draw[vector] (0,0)--(2.3,0) node[right] {$\vect v$};
   \draw[vector] (0,0)--(-\chi:2.3) node[below right] {$\vect v'$};
   \draw[->] (.8,0) arc(0:-\chi:.8);
   \node at (-.5*\chi:1.14) {$\theta$};
   \draw[guide] (-\chi:2.3)--({2.3*cos(\chi)},0);
   \node[above] at (1.0,.20) {$v'=v$};
   \node at (1.15,-2.45) {纵向变化 $v-v\cos\theta$};
  \end{scope}
\end{tikzpicture}
\figcap{引力双曲线与远处速度夹角。右图把远处速度平移到同一起点比较，两者大小相等。}
\label{ch11:fig:scattering}
\end{minipage}\end{center}

由于 $b=a_g\cot(\theta/2)$、$\dd\Omega=2\pi\sin\theta\,\dd\theta$，轴对称散射给出
\begin{align*}
 \frac{\dd\sigma}{\dd\Omega}
 &=\frac{b}{\sin\theta}\left|\frac{\dd b}{\dd\theta}\right|\\
 &=\boxed{\left(\frac{GM_s}{2v^2}\right)^2
 \csc^4\frac\theta2}.
\end{align*}
这里使用了 $|\dd b/\dd\theta|=(a_g/2)\csc^2(\theta/2)$ 以及
$\sin\theta=2\sin(\theta/2)\cos(\theta/2)$。尘埃质量在速度与能量关系中约去，
因此截面与单个尘粒质量无关；结果量纲为面积。

几何半径与远处碰撞截面半径不同。恰好碰到表面的轨道在最近处速度为 $b_{\rm hit}v/r_s$，
由能量守恒
\[
 \frac{v^2}{2}=\frac{b_{\rm hit}^2v^2}{2r_s^2}-\frac{GM_s}{r_s}
 \quad\Longrightarrow\quad b_{\rm hit}^2=r_s^2+2a_gr_s.
\]
所以直接撞击的质量流率为 $\rho v\pi b_{\rm hit}^2$。在“吸收并停到飞船参考系”的假设下，每单位质量损失纵向动量 $v$，故
\[
 F_{\rm hit}=\pi b_{\rm hit}^2\rho v^2
 =\pi\rho v^2(r_s^2+2a_gr_s)\simeq\boxed{\pi r_s^2\rho v^2}.
\]
末一步与题设的小角近似一致，要求 $a_g\ll r_s$，在这一阶忽略引力聚焦对直接碰撞的修正。
冲量参数在 $b$ 到 $b+\dd b$ 的掠射尘埃，单位时间通过的质量为
$\rho v\,2\pi b\dd b$，纵向速度损失为 $v(1-\cos\theta)$。又
\[
 1-\cos\theta=\frac{2a_g^2}{b^2+a_g^2},
\]
所以排除直接碰撞区域后
\begin{align*}
 F_{\rm grav}
 &=\rho v^2\int_{b_{\rm hit}}^{R_d}2\pi b(1-\cos\theta)\,\dd b\\
 &=2\pi\rho v^2a_g^2
 \ln\!\frac{R_d^2+a_g^2}{b_{\rm hit}^2+a_g^2}.
\end{align*}
当所有散射角都很小，即 $a_g\ll r_s<R_d$ 时，化为
\[
 \boxed{F_{\rm grav}\simeq
 \frac{4\pi\rho G^2M_s^2}{v^2}\ln\frac{R_d}{r_s}}.
\]
总阻力为 $F_{\rm hit}+F_{\rm grav}$，方向与飞船相对尘埃的速度相反。若尘埃发生镜面反弹、烧蚀或强引力聚焦，
直接碰撞项的系数和积分下限都需相应修正；上述精确积分要求 $R_d>b_{\rm hit}$，小角结果还用到 $b_{\rm hit}\simeq r_s$。
碰撞区不能再计入掠射阻力，否则会重复计算。
\end{xesolution}

\Needspace{16\baselineskip}\begin{origexamplebox}{弱线性径向附加力造成的近日点进动}
\label{ch11:ex:07}
除太阳质量 $M$ 的平方反比引力外，质量为 $m$ 的行星还受到很小的径向外力
$\vect F_1=kmr\,\uvec r$，其中 $k>0$ 为常量。设近圆轨道半径为 $R$，求近日点进动角速度。

\ThQPrecession
\end{origexamplebox}
\Needspace{6\baselineskip}\begin{xesolution}
\ThAPrecession

附加力向外，故单位质量势能可写为
\[
 \Phi(r)=-\frac{GM}{r}-\frac12kr^2.
\]
在半径 $R$ 的圆轨道上，向心条件给出实际平均角频率
\[
 \boxed{\Omega^2=\frac{GM}{R^3}-k}.
\]
为求近日点进动，必须把径向小振动频率与 $\Omega$ 区分。令单位质量角动量为 $h=R^2\Omega$，
有效势为 $\Phi_{\rm eff}=\Phi+h^2/(2r^2)$。逐次求导得
\[
 \Phi_{\rm eff}'=\frac{GM}{r^2}-kr-\frac{h^2}{r^3},\qquad
 \Phi_{\rm eff}''=-\frac{2GM}{r^3}-k+\frac{3h^2}{r^4}.
\]
圆轨道条件使 $\Phi_{\rm eff}'(R)=0$ 且 $h^2=GMR-kR^4$。
写 $r=R+x$，将 $\ddot r=-\Phi_{\rm eff}'(r)$ 保留到 $x$ 的一阶，得到
\[
 \ddot x+\kappa_r^2x=0,
 \qquad
 \kappa_r^2=\Phi_{\rm eff}''(R)
 =\frac{GM}{R^3}-4k.
\]
因此一次径向周期内方位角略多于 $2\pi$，近日点沿运动方向进动，其角速度为
\[
 \boxed{\dot\varpi=\Omega-\kappa_r
 =\sqrt{\frac{GM}{R^3}-k}
 -\sqrt{\frac{GM}{R^3}-4k}}.
\]
记无扰动平均运动 $n=\sqrt{GM/R^3}$。在 $k\ll n^2$ 时，
$\Omega=n-k/(2n)+O(k^2/n^3)$，
$\kappa_r=n-2k/n+O(k^2/n^3)$，相减得
\[
 \boxed{\dot\varpi\simeq\frac{3k}{2n}
 =\frac{3k}{2}\sqrt{\frac{R^3}{GM}}}.
\]
稳定近圆运动还要求 $GM/R^3>4k$。当 $k\to0$ 时 $\Omega=\kappa_r$、进动消失，极限正确。
\end{xesolution}

\Needspace{16\baselineskip}\begin{origexamplebox}{弱速度阻力下椭圆轨道的长期演化}
\label{ch11:ex:08}
质量为 $m$ 的质点在固定质量 $M$ 的平方反比引力下沿椭圆运动，同时受到很小的阻力
$\vect f=-\gamma\vect v$。取时间间隔 $\dd t$ 远长于轨道周期、又短到轨道根数的相对变化仍很小。
已知机械能和角动量的大小分别为 $E<0,L$。求周期平均后的 $\dd E,\dd L$，并证明轨道只按比例缩小、偏心率不变。

\ThQOrbitDrag
\end{origexamplebox}
\Needspace{6\baselineskip}\begin{xesolution}
\ThAOrbitDrag

阻力的瞬时功率为
\[
 \dot E=\vect f\cdot\vect v=-\gamma v^2.
\]
一周内先忽略微小阻力对轨道形状的改变。沿这一闭合开普勒轨道，
\[
 \frac{\dd}{\dd t}(\vect r\cdot\vect v)
 =v^2+\vect r\cdot\dot{\vect v}=v^2-\frac{GM}{r}.
\]
一周后 $\vect r\cdot\vect v$ 回到原值，故左端周期平均为零，
$\langle v^2\rangle=GM\langle1/r\rangle$。代入
$E=m\langle v^2\rangle/2-GMm\langle1/r\rangle$，
得 $E=-m\langle v^2\rangle/2$。于是长期平均为
\[
 \boxed{\dd E=\frac{2\gamma}{m}E\,\dd t}.
\]
因为 $E<0$，右端为负，轨道能量变得更负。阻力矩则直接给出
\[
 \dot{\vect L}=\vect r\times(-\gamma\vect v)
 =-\frac\gamma m\vect L,
 \qquad
 \boxed{\dd L=-\frac\gamma mL\,\dd t}.
\]

开普勒椭圆的偏心率满足
\[
 e^2=1+\frac{2EL^2}{G^2M^2m^3}.
\]
在当前精度下
\[
 \frac{\dd(EL^2)}{EL^2}
 =\frac{\dd E}{E}+2\frac{\dd L}{L}
 =\frac{2\gamma}{m}\dd t-\frac{2\gamma}{m}\dd t=0,
\]
所以
\[
 \boxed{\dd e=0}.
\]
另一方面 $E=-GMm/(2a)$，从而
\[
 \frac{\dd a}{a}=-\frac{\dd E}{E}
 =-\frac{2\gamma}{m}\dd t.
\]
半长轴按比例减小而偏心率不变，因此近、远点距离 $a(1\mp e)$ 以相同比例收缩。这个结论是周期平均结果，
单个轨道周期内的瞬时轨道根数仍会有小幅振荡。

把半长轴方程积分，得到
\[
 \boxed{a(t)=a_0\exp\!\left(-\frac{2\gamma t}{m}\right)}.
\]
由开普勒第三定律 $T\propto a^{3/2}$，轨道周期随时间按
\[
 \boxed{T(t)=T_0\exp\!\left(-\frac{3\gamma t}{m}\right)}
\]
缩短。这个指数规律并不意味着质点在有限时间撞到中心；在该平均模型中 $a$ 只在
$t\to\infty$ 时趋于零。

还可直接核对能量损失率：阻力功率为
$\dot E=\vect F_d\cdot\vect v=-\gamma v^2$。椭圆轨道一周平均满足
$\langle v^2\rangle=GM/a$，于是
\[
 \langle\dot E\rangle=-\frac{\gamma GM}{a},
\]
与 $E=-GMm/(2a)$ 及 $\dot a/a=-2\gamma/m$ 完全一致。
上述推导要求一周内轨道根数变化很小，即 $\gamma T/m\ll1$；阻尼过强时不能再把运动看成一串缓慢收缩的开普勒椭圆。
\end{xesolution}

\endgroup

\hypersetup{bookmarksdepth=2}
